%==============================================================================
% Chapter 2 -- Computational design of protein nanoparticle vaccines
%==============================================================================
\chapter{Computational design of protein nanoparticle vaccines}
\label{ch:comp-design}

% -----------------------------------------------------------------------------
% This chapter expands the "Computational design of protein nanoparticles"
% thread of the Design-CP paper's Related Work into a proper background section,
% and then states the objective of the project. It is the pivot of the report:
% it sets up why modelling the *whole assembly* end to end is the enabling first
% step, which is exactly the premise of the Design-CP work reported in
% Chapters 3-4.
% -----------------------------------------------------------------------------

\section{Current methods}
\label{sec:current-methods}

State-of-the-art design of protein nanoparticles follows a two-stage
``dock-and-design'' recipe inherited from the Rosetta era. First, symmetric
oligomeric building blocks are rigidly docked into a target point-group
architecture (for example a tetrahedron, octahedron or icosahedron) so that their
symmetry axes coincide with those of the target; second, the newly created
subunit--subunit interfaces are redesigned to drive stable self-assembly. This
programme produced the first self-assembling protein nanomaterials with
atomic-level accuracy \parencite{king_computational_2012}, its multi-component
extension \parencite{king_accurate_2014}, megadalton-scale two-component
icosahedral complexes such as I53-50 \parencite{bale_accurate_2016} and the
hyperstable 60-subunit icosahedron I3-01 \parencite{hsia_design_2016}, building on
symmetric self-assembly principles established earlier
\parencite{padilla_nanohedra_2001}. Machine learning has since accelerated the
generation of building blocks and the design of interfaces: symmetric assemblies
have been produced by hallucination
\parencite{wickyHallucinatingSymmetricProtein2022}, interfaces are now routinely
designed with ProteinMPNN \parencite{de_haas_rapid_2024}, and the building blocks
themselves are drawn from AlphaFold2-predicted structures
\parencite{haas_sequence_2025} or generated de novo with RFdiffusion
\parencite{haas_novo_2026}. Related advances extend the repertoire to
pseudosymmetric and multi-component cages
\parencite{dowlingHierarchicalDesignPseudosymmetric2025,
rankovicComputationalDesignBifaceted2025,lee_four-component_2025,kibler_design_2024}.

Despite this progress, essentially all of these pipelines share a structural
limitation: they still rely on a \emph{rigid-body docking step} over a library of
pre-computed oligomers. The building blocks are designed, or predicted, in
isolation and are only afterwards assembled by a search over rigid poses. The
inter-subunit interfaces---the very interactions that determine whether, and how
well, a particle assembles---are therefore never generated jointly with the rest
of the structure. A natural next step is to remove the docking stage altogether
and model the entire assembly, including all of its inter-subunit interfaces, with
a single generative network.


\section{Objective of our project}
\label{sec:objective}

The objective of this project is to marry the fully protein-based nanoparticle
vaccine platform of Chapter~\ref{ch:introduction} with the generative-modelling
capabilities of Chapter~\ref{sec:protein-design}: to design protein nanoparticle
vaccines \emph{end to end}, as whole assemblies, with a single all-atom generative
model rather than a docking pipeline over independently designed parts. In the
longer term (Chapter~\ref{ch:future-work}) this means \emph{conditional} design---
scaffolding a chosen antigen onto a self-assembling particle---but that ambition
rests on a prerequisite that turns out to be non-trivial.

\paragraph{The enabling first step: modelling the whole assembly.}
Before one can condition generation on a displayed antigen, one must first be able
to \emph{model the whole assembly at all}. All-atom generative models can, in
principle, design a large multimeric complex by jointly denoising every chain and
every atom; in practice, their quadratic token-pair and atom-pair representations
exhaust the memory of a single GPU as the number of subunits and residues grows.
A full icosahedral particle---sixty chains and many thousands of residues---simply
does not fit. This memory ceiling is the reason the field still falls back on
docking pre-computed oligomers rather than generating the assembly directly.

This report addresses exactly that bottleneck. Building on the all-atom model
RFdiffusion3 \parencite{butcher_novo_2025}, we borrow \emph{context parallelism}
from large-language-model and structure-prediction systems
\parencite{lin_fold-cp_2026,dao_flashattention_2022} to distribute the quadratic
activations across a multi-GPU mesh, so that the full assembly can be denoised in a
single trajectory. The result---presented in Chapters~\ref{ch:methods}
and~\ref{ch:results}---is, to our knowledge, the first end-to-end, all-atom
generative design of symmetric protein assemblies that models the full set of
inter-subunit interactions without an intermediate docking step. Establishing that
we can reliably model the whole assembly is the premise on which every subsequent
research direction in this report is built.
